\section{Related Work}\label{sec:related}

\paragraph{Instance categorization for imbalanced learning.}
Several frameworks decompose training data by learning difficulty.
Instance Hardness \cite{smith2014instance} assigns a difficulty score per
instance from majority votes across multiple learners. Dataset Cartography
\cite{swayamdipta2020dataset} maps instances by training-dynamics
confidence/variability. Data-IQ \cite{seedat2022dataiq} extends this to
tabular data via prediction trajectories. The closest line of work is
Napiera{\l}a \& Stefanowski \cite{napierala2010learning,napierala2016types},
who categorize minority instances into safe, borderline, rare, and outlier
based on the proportion of same-class neighbours in a k-nearest-neighbour
graph; S\'aez et al.\ \cite{saez2016analyzing} extend this to multi-class
imbalanced data and analyze how oversampling interacts with each type. Crucially,
all of these characterize instance \emph{difficulty}; none condition on the model
being \emph{wrong} and ask \emph{which remedy applies}. Our decomposition differs
in three ways: (i) it operates on the \emph{misclassified minority} and maps each
error to a remedy --- move the threshold, add data, or accept it --- rather than
scoring difficulty; (ii) it introduces a \emph{threshold-recoverable} class (the
model already ranks the instance correctly and only the default cutoff is wrong),
an axis absent from prior difficulty taxonomies; and (iii) it combines a
model-based aleatoric--epistemic decomposition --- separating Bayes-rate ambiguity
from sample-size limits, which k-NN typing cannot --- with a local class-ratio
measure, used as a diagnostic rather than as a competing oversampler.

\paragraph{Threshold-moving, cost-sensitivity, and the limits of resampling.}
A long line of work argues the decision threshold, not the training distribution,
is the right lever under imbalance. Provost \cite{provost2000imbalanced} warns
against the default $0.5$ cutoff, and Elkan \cite{elkan2001foundations} shows that
changing class proportions is \emph{equivalent} to changing the decision
threshold/costs, so rebalancing barely alters the learned scoring function;
Drummond and Holte \cite{drummond2003c45} reach the analogous conclusion via cost
curves. Large-scale and clinical studies confirm the empirical consequence:
resampling does not improve ranking/AUC and its apparent gains are reproduced by a
threshold shift, while it degrades probability calibration
\cite{elor2022smote,goorbergh2022harm,dalpozzolo2015calibrating,carriero2025harms}.
We do \emph{not} claim these regularities as new --- we reproduce them across our
roster and instead \emph{explain} them: our reducibility decomposition shows the
misclassified minority is predominantly \emph{threshold-recoverable}, which is
\emph{why} a threshold move suffices and oversampling behaves as an implicit,
inefficient enactment of it. Our contribution is the per-instance instrument that
accounts for this equivalence and locates its exceptions (the data-reducible and
irreducible minority), not the equivalence itself.

\paragraph{Class overlap and imbalance.}
Santos et al.\ \cite{santos2023unifying} provide a unifying multi-view
panorama of the relationship between class overlap and class imbalance,
arguing that overlap is often the underlying driver of imbalance-related
difficulty. Their framing supports our boundary-generation diagnosis: when
SMOTE generates near class-overlap regions, the harm to accuracy is a
predictable consequence of overlap-blind seeding.

\paragraph{Characteristic-driven and complexity-guided selection.}
Komorniczak, Ksieniewicz, and Wo{\'z}niak \cite{komorniczak2022data}
demonstrate that data-complexity metrics correlate with the accuracy gain
oversampling delivers, and Carvalho et al.\ \cite{carvalho2025resampling} survey
how resampler performance tracks data structure. Most directly, Jiang and Ye
\cite{jiang2026beyond} control the imbalance ratio and show data characteristics
moderate oversampler choice --- yet their resulting selection rule does \emph{not}
beat always-SMOTE, and they leave open whether a held-out, characteristic-driven
selector can. This line is \emph{correlational}: it relates measured complexity to
observed performance without explaining the mechanism. Our contribution is
orthogonal --- a mechanism-level account of \emph{why} boundary generation hurts --- and we
answer the open selection question directly, on a full strategy menu with held-out
validation (\S\ref{sec:results:selection}), finding that the menu, not the
selector, is the dominant lever. Our account is mechanism-level and
intervention-supported rather than correlational.

\paragraph{SMOTE variants and cleaning.}
Borderline-SMOTE \cite{han2005borderline} restricts generation to instances
near the decision boundary; ADASYN \cite{he2008adasyn} upweights generation
in sparse minority regions. Both methods \emph{increase} generation in
exactly the regions we identify as harmful. Safe-Level-SMOTE
\cite{bunkhumpornpat2009safe} uses each seed's ``safe level'' (count of
same-class neighbours) to gate interpolation, an idea closely related to
ours but instantiated via k-NN heuristics rather than uncertainty
decomposition; we compare directly in \S\ref{sec:results}.
Post-hoc cleaning methods combine SMOTE with instance removal: SMOTE+ENN
\cite{batista2004study,wilson1972asymptotic} removes instances misclassified
by edited nearest neighbours, while SMOTE+Tomek
\cite{batista2004study,tomek1976two} removes Tomek-link pairs. These methods
are agnostic to \emph{why} an instance is problematic, removing both
correctable and boundary instances indiscriminately. Beyond the canonical
family, MWMOTE \cite{barua2014mwmote}, ProWSyn \cite{barua2011prowsyn}, and
polynomial-fitting SMOTE \cite{gazzah2008new} represent more recent
extensions. The 85-variant survey of Kov\'acs \cite{kovacs2019smote}
provides a broader perspective; we include the top three of his ranked
variants as baselines. Theoretically, Sakho et al.\ \cite{sakho2024rebalancing}
show that \emph{default} SMOTE asymptotically copies minority points and that its
synthetic density \emph{vanishes} near the boundary of the minority support; our
boundary-crossing account is therefore specific to the \emph{overlap} regime ---
interpolation between minority seeds that already lie in class-overlap territory,
where the segment crosses into majority-posterior space
(Proposition~\ref{prop:cross}) --- and is complementary to, not in conflict with, their
clean-margin analysis.

\paragraph{Label noise detection.}
Confident Learning \cite{northcutt2021confident} estimates the joint
distribution of noisy and true labels to identify mislabeled instances. Our
Cat\,1 (noise) category captures a similar signal but is embedded in a
broader three-way decomposition that distinguishes noise from boundary
ambiguity --- a distinction Confident Learning does not make.

\paragraph{Uncertainty decomposition.}
Shaker and H\"ullermeier \cite{shaker2020aleatoric} decompose random-forest
predictive uncertainty into aleatoric and epistemic components via
ensemble-of-trees disagreement. H\"ullermeier and Waegeman
\cite{hullermeier2021aleatoric} provide a general treatment of the
aleatoric--epistemic distinction, and Nguyen et al.\
\cite{nguyen2019epistemic} use it for data acquisition --- epistemic uncertainty
marks the \emph{reducible} part an extra label can resolve, aleatoric the
irreducible part --- in an active-learning setting. We adapt this
reducible-versus-irreducible logic to imbalanced learning: combined with a local
class-ratio measure and a threshold-recoverability test, it yields the
remedy-mapped decomposition of the misclassified minority that is our headline
contribution.

\paragraph{Recent reviews.}
Fern\'andez et al.\ \cite{fernandez2018smote} mark the 15-year anniversary
of SMOTE with a progress-and-challenges review. Carvalho et al.\
\cite{carvalho2025resampling} survey resampling approaches from a data
perspective; their review reinforces the observation that resampler
performance is best predicted by data structure rather than by the
algorithmic family of the resampler --- a perspective consistent with our
exclusion-beats-targeting finding.
